% Pista ana-26s-q1 · Anàlisi
% Procedència: PAU Catalunya 2026 · Convocatòria de setembre · Sèrie 2 · Exercici 1
\documentclass[11pt,a4paper]{article}
\usepackage{pista}
\pistainfo{Catalunya 2026}{Convocatòria de setembre}{Sèrie 2}

\begin{document}

\section*{\textcolor{blauPista}{Exercici 1}}

\noindent Considereu la funció $f(x) = \dfrac{5x}{x^{2} + 1}$.

\subsection*{1.1. \normalfont Estudieu-ne el domini i els intervals de creixement i decreixement, així com els possibles màxims i mínims locals. \hspace*{\fill}\textnormal{[1~punt]}}

\pista{Segueix l'ordre de l'estudi d'una funció. Aquí et demanen dos passos:
\begin{itemize}
\item \textit{Domini.} Mira si el denominador $x^{2} + 1$ es pot anul·lar per a algun valor real.
\item \textit{Creixement, decreixement i extrems relatius.} Deriva amb la regla del quocient i simplifica el numerador amb cura (és on és més fàcil equivocar-se de signe). El denominador d'$f'(x)$ és un quadrat i no s'anul·la: el signe d'$f'$ depèn només del numerador. Troba'n els zeros i fes una taula de signes.
\end{itemize}

\medskip
\textbf{Atenció:} dona també l'ordenada dels extrems, calculant-la a $f(x)$ (no a $f'(x)$).}

\subsection*{1.2. \normalfont Calculeu les equacions de totes les asímptotes d'aquesta funció (si en té) i dibuixeu un esbós de la seva gràfica. \hspace*{\fill}\textnormal{[0,75~punts]}}

\pista{Tens gairebé tota la feina feta a l'apartat 1.1. Et falten les asímptotes:
\begin{itemize}
\item \textit{Verticals.} Creus que el denominador s'anul·la per algun valor? Només en aquests valors hi pot haver una asímptota vertical.
\item \textit{Horitzontals.} Prova valors grans: $f(10) \approx 0{,}5$ i $f(100) \approx 0{,}05$. Cap a quin valor s'acosta $f$? I amb $x = -100$? Després, justifica-ho amb el límit.
\end{itemize}
Per a l'esbós, situa el màxim i el mínim de l'apartat 1.1 i fes que la corba s'acosti a l'asímptota per tots dos costats. No et cal la segona derivada.}

\subsection*{1.3. \normalfont Calculeu l'equació de la recta normal (és a dir, la recta perpendicular a la tangent) a la corba $y = f(x)$ en el punt d'abscissa $x = 2$. \hspace*{\fill}\textnormal{[0,75~punts]}}

\pista{La normal és la recta perpendicular a la tangent, amb pendent $m = -\dfrac{1}{f'(2)}$:
\[
y - f(2) = m\,(x - 2)
\]
Calcula el punt $\bigl(2, f(2)\bigr)$ i el pendent de la tangent, $f'(2)$, amb la derivada de l'apartat 1.1.

\medskip
\textbf{Atenció:} cal invertir el valor \emph{i} canviar-li el signe. L'error típic és fer només una de les dues coses, o fer servir directament $f'(2)$.}

\end{document}
